\documentclass[11pt, a4paper]{article}

\usepackage[top = 1 in, bottom = 1 in, left = 0.5 in, right = 0.5 in]{geometry}

\usepackage{amsmath, amssymb, amsfonts}
\allowdisplaybreaks[4]

\usepackage{enumerate}
\usepackage{multirow}
\usepackage{hhline}
\usepackage{array}
\usepackage{longtable}
\usepackage{graphicx}
\usepackage{tabularray}
\usepackage{undertilde}
\usepackage{dingbat}
\usepackage{fontawesome5}
\usepackage[colorlinks=true, linkcolor=blue, urlcolor=red]{hyperref}
\usepackage{tasks}
\usepackage{bbding}
\usepackage{twemojis}
% how to use bull's eye ----- \scalebox{2.0}{\twemoji{bullseye}}
\usepackage{customdice}
% how to put dice face ------ \dice{2}
\usepackage{fontspec}
\usepackage{fancyhdr}
\usepackage{lastpage}
\usepackage{skak}
\usepackage{mathrsfs}
\usepackage{xcolor}

\pagestyle{fancy}
\fancyhf{}
\fancyfoot[C]{Page \thepage\ of \pageref{LastPage}}
\renewcommand{\headrulewidth}{0pt}



\title{Expectation-Maximization Algorithm}
\author{{\fontsize{25}{20} \benyorithafont Ananda Biswas}}
\date{}

\newfontface\gloriahallelujahfont{GLORIAHALLELUJAH-REGULAR.TTF}

\newfontface\benyorithafont{Benyoritha-G334D.otf}


\begin{document}

\maketitle

\section*{\faArrowAltCircleRight[regular] \textcolor{blue}{Censored Exponentially Distributed Survival Times}}

We suppose $W$ is a non-negative random variable having an exponential
distribution with \colorbox{yellow}{mean $\mu$}. Thus its probability density function
(P.D.F.) is given by

\begin{equation} \label{Exp_PDF}
f(w;\mu) = \dfrac{1}{\mu} \exp\!\left\{-\dfrac{w}{\mu}\right\} I_{(0,\infty)}(w), \qquad (\mu>0),
\end{equation}

where the indicator function

\[
I_{(0,\infty)}(w)
=
\begin{cases}
1, & \text{for } w>0,\\
0, & \text{elsewhere.}
\end{cases}
\]

The distribution function is given by

\[
F(w;\mu)
=
\left\{
1-\exp\!\left(-\dfrac{w}{\mu}\right)
\right\}
I_{(0,\infty)}(w).
\]

In survival or reliability analyses, a study to observe a random sample
$W_1,\ldots,W_n$ from (\ref{Exp_PDF}) will generally be terminated in practice before all of these random variables are able to be observed. We let

\[
\boldsymbol{y}
=
\left(
\boldsymbol{y}_1^{{T}},
\ldots,
\boldsymbol{y}_n^{{T}}
\right)^{{T}}
\]

denote the observed data, where

\[
\boldsymbol{y}_j
=
(c_j,\delta_j)^{{T}}
\]

and $\delta_j = 0$ or $1$ according as the observation $W_j$ is censored
or uncensored at $c_j$ $(j=1,\ldots,n)$. That is, if the observation
$W_j$ is uncensored, its realized value $w_j$ is equal to
$c_j$, whereas if it is censored at $c_j$, then $w_j$ is some value
greater than $c_j$ $(j=1,\ldots,n)$.

We suppose now that the observations have been relabelled so that
$W_1,\ldots,W_r$ denote the $r$ uncensored observations and
$W_{r+1},\ldots,W_n$ the $n-r$ censored observations. The likelihood function for $\mu$ formed on the basis of $\boldsymbol{y}$ is given by

\begin{align}
L(\mu)
&=
\prod \limits_{j = 1}^{r} f(w_j; \mu) \cdot \prod \limits_{j = r+1}^{n} \left( 1 - F(c_j; \mu) \right) \nonumber \\[0.5em]
&=
\prod \limits_{j = 1}^{r} \dfrac{1}{\mu} \exp\!\left\{-\dfrac{w_j}{\mu}\right\} \cdot \prod \limits_{j = r+1}^{n} \exp\!\left\{-\dfrac{c_j}{\mu}\right\} \nonumber \\[0.5em]
&=
\left( \dfrac{1}{\mu} \right)^r \cdot \exp\!\left\{- \sum \limits_{j = 1}^{n} \dfrac{c_j}{\mu} \right\}
\end{align}

\newpage

The \colorbox{yellow}{log likelihood} function for $\mu$ formed on the basis of
$\boldsymbol{y}$ is given by

\begin{equation} \label{loglikelihood}
\ell(\mu)
=
\log L(\mu)
=
-r \log \mu
-
\sum \limits_{j=1}^{n} \dfrac{c_j}{\mu}.
\end{equation}

In this case, the MLE of $\mu$ can be derived explicitly from
equating the derivative of (\ref{loglikelihood}) to zero to give

\[
\hat{\mu}
=
\sum \limits_{j=1}^{n} \dfrac{c_j}{r}.
\]

Thus there is no need for the iterative computation of $\hat{\mu}$.
But in this simple case, it is instructive to demonstrate how the
EM algorithm would work.

The complete-data vector $\boldsymbol{x}$ can be declared to be

\begin{align*}
\boldsymbol{x}
&=
(w_1,\ldots,w_n)^T \\[0.5em]
&=
(w_1,\ldots,w_r,\boldsymbol{z}^T)^T,
\end{align*}

where

\[
\boldsymbol{z}
=
(w_{r+1},\ldots,w_n)^T
\]

contains the unobservable realizations of the $n-r$ censored random
variables. In this example, the so-called unobservable or missing
vector $\boldsymbol{z}$ is potentially observable in a data sense,
as if the experiment were continued until each item failed, then there
would be no censored observations.

The \colorbox{yellow}{complete-data log likelihood} is given by

\begin{align*}
\ell_c(\mu)
=
\log L_c(\mu)
&=
\sum \limits_{j=1}^{n}
\log f(w_j;\mu) \\[0.5em]
&=
-n \log \mu
-
\dfrac{1}{\mu}
\sum \limits_{j=1}^{n} w_j .
\end{align*}

As $\ell_c(\mu)$ can be seen to be linear in the unobservable data
$w_{r+1},\ldots,w_n$, the calculation of
$Q(\mu;\mu^{(k)})$ on the E-step (on the $(k+1)$th iteration)
simply requires each such $w_j$ to be replaced by its conditional
expectation given the observed data $\boldsymbol{y}$, using the
current fit $\mu^{(k)}$ for $\mu$ (Why ?).

By the lack of memory of the exponential distribution,
the conditional distribution of $W_j-c_j$ given that $W_j>c_j$
is still exponential with mean $\mu$. Then, we have that

\begin{align}
\colorbox{yellow}{$E_{\mu^{(k)}}(W_j \mid \boldsymbol{y})$}
&=
E_{\mu^{(k)}}(W_j \mid W_j>c_j) \nonumber \\[0.5em]
&=
c_j + E_{\mu^{(k)}}(W_j) \nonumber \\[0.5em]
&=
c_j + \mu^{(k)} \label{conditional_expectation}
\end{align}

for $j=r+1,\ldots,n$.

On using (\ref{conditional_expectation}) to take the current conditional expectation of the complete-data log likelihood $\ell_c(\mu)$, we have that

\begin{align}
\colorbox{yellow}{$Q(\mu;\mu^{(k)})$}
&=
-n \log \mu
-
\dfrac{1}{\mu}
\left\{
\sum \limits_{j=1}^{r} c_j
+
\sum \limits_{j=r+1}^{n}
\left(c_j+\mu^{(k)}\right)
\right\}  \nonumber \\[0.5em]
&=
-n \log \mu
-
\dfrac{1}{\mu}
\left\{
\sum \limits_{j=1}^{n} c_j
+
(n-r)\mu^{(k)}
\right\}. \label{Lower_Bound}
\end{align}

Concerning the M-step on the $(k+1)$th iteration, it follows from
(\ref{Lower_Bound}) that the value of $\mu$ that maximizes
$Q(\mu;\mu^{(k)})$ is given by the MLE of $\mu$ that would be formed
from the complete data, but with each unobservable $w_j$ replaced by its
current conditional expectation given by (\ref{conditional_expectation}). Accordingly,

\begin{align*}
\colorbox{yellow}{$\mu^{(k+1)}$}
&=
\left\{
\sum \limits_{j=1}^{r} c_j
+
\sum \limits_{j=r+1}^{n}
E_{\mu^{(k)}}(W_j \mid \boldsymbol{y})
\right\}
\Bigg/n
\\[1.2ex]
&=
\left\{
\sum \limits_{j=1}^{r} c_j
+
\sum \limits_{j=r+1}^{n}
\left(c_j+\mu^{(k)}\right)
\right\}
\Bigg/n
\\[1.2ex]
&=
\left\{
\sum \limits_{j=1}^{n} c_j
+
(n-r)\mu^{(k)}
\right\}
\Bigg/n.
\end{align*}

\end{document}